%% kapa: exam=1 round=37 year=2026 period=2 form=A subject=accounting subject_ko=회계학 seq=9 number=49
%% source: Q-Net 감정평가사 기출문제 (cst00309) artlSeq=5258924 file=raw/1008/5258924/2026년 제37회 감정평가사 1차 2교시.pdf
%% conversion: kordoc 4.24.51 -> tools/gen1.py (auto); review=proofread-sonnet; flags=none
\begin{question}{49}
무형자산에 관한 설명으로 옳은 것은?
\begin{choices}
  \item 내용연수가 유한한 무형자산의 잔존가치는 해당 자산의 장부금액과 같을 수는 있으나, 장부금액보다 더 클 수는 없다.
  \item 내부적으로 창출한 영업권도 무형자산으로 인식할 수 있다.
  \item 내용연수가 비한정인 무형자산을 유한 내용연수로 재평가하는 경우에는 자산손상의 징후에 해당되지 않으므로 손상차손을 인식하지 않는다.
  \item 경제적 효익이 소비될 것으로 예상되는 형태를 신뢰성 있게 결정할 수 없는 내용연수가 비한정인 무형자산은 정액법을 적용하여 상각한다.
  \item 내용연수가 유한한 무형자산을 내용연수 종료 시점에 제3자가 구입하기로 약정한 경우, 잔존가치는 영(0)으로 보지 않는다.
\end{choices}
\end{question}
\answer{5}
